\begin{tcolorbox}[colback=cadmiumgreen!5!white,colframe=forestgreen!75!white,breakable,title=\textit{Basic setting prompt template of optimized with in-context for the mortality prediction task on the MIMIC-IV.}]
\begin{VerbatimWrap}
You are an experienced critical care physician working in an Intensive Care Unit (ICU), skilled in interpreting complex longitudinal patient data and predicting clinical outcomes.
System Prompt: You are an experienced critical care physician working in an Intensive Care Unit (ICU), skilled in interpreting complex longitudinal patient data and predicting clinical outcomes.

User Prompt: I will provide you with longitudinal medical information for a patient. The data covers 3 visits that occurred at 2113-01-31, 2113-02-01, 2113-02-02.
Each clinical feature is presented as a list of values, corresponding to these visits. Missing values are represented as `NaN` for numerical values and "unknown" for categorical values. Note that units and reference ranges are provided alongside relevant features.

Patient Background:
- Sex: male
- Age: 50 years

Your Task:
Your primary task is to assess the provided medical data and analyze the health records from ICU visits to determine the likelihood of the patient not surviving their hospital stay.

Instructions & Output Format:
Please first perform a step-by-step analysis of the patient data, considering trends, abnormal values relative to reference ranges, and their clinical significance for survival. Then, provide a final assessment of the likelihood of not surviving the hospital stay.

Your final output must be a JSON object containing two keys:
1.  `"think"`: A string containing your detailed step-by-step clinical reasoning (under 500 words).
2.  `"answer"`: A floating-point number between 0 and 1 representing the predicted probability of mortality (higher value means higher likelihood of death).

Example Format: ```json { "think": "The patient presents with worsening X, stable Y, and improved Z. Factor A is a major risk indicator... Overall assessment suggests a high risk.", "answer": 0.85 }```

Handling Uncertainty:
In situations where the provided data is clearly insufficient or too ambiguous to make a reasonable prediction, respond with the exact phrase: `I do not know`.

Now, please analyze and predict for the following patient:

Clinical Features Over Time:
- Capillary refill rate: [0.0, 0.0, 0.0]
- Glascow coma scale eye opening: [Spontaneously, Spontaneously, To Speech]
- Glascow coma scale motor response: [Obeys Commands, Obeys Commands, Obeys Commands]
- Glascow coma scale total: [0.0, 0.0, 0.0]
- Glascow coma scale verbal response: [Oriented, Oriented, No Response]
- Diastolic blood pressure: [55.0, 74.0, 73.0]
- Fraction inspired oxygen: [80.0, 50.0, 70.0]
- Glucose: [119.0, 118.0, 127.0]
- Heart Rate: [86.0, 110.0, 118.0]
- Height: [157.0, NaN, NaN]
- Mean blood pressure: [67.0, 85.0, 102.0]
- Oxygen saturation: [96.0, 100.0, 100.0]
- Respiratory rate: [17.0, 26.0, 15.0]
- Systolic blood pressure: [105.0, 124.0, 156.0]
- Temperature: [37.89, 37.61, 37.17]
- Weight: [80.92, 80.92, 80.92]
- pH: [7.48, 7.47, 7.51]
\end{VerbatimWrap}
\end{tcolorbox}




